\subsection*{A countable list for reversible classes generated by polyGEMs}

\begin{claim}{D.3.7}{Proved}{Abelian groups do not retain arbitrary-cardinal distinctions}
\label{D37GEMnormalform}
Put $C_{p^\infty}=\mathbb Z[1/p]/\mathbb Z$.  There is a countable
list of countable connected spaces
\[
 \mathcal E=\{B\mathbb Q\}\ \cup
 \{K(\mathbb Z_{(p)},n),\ K(\mathbb F_p,n),\
                  K(C_{p^\infty},n):p\text{ prime},\ n\geq1\}
 \tag{D.3.7a}
\]
with the following property.  For every abelian group $A$ and $n\geq1$,
the smallest reversible fundamental localizer containing
$K(A,n)\to *$ is generated reversibly by a subset of
\[
 \{B\mathbb Q\}\cup
 \{K(\mathbb Z_{(p)},n),K(\mathbb F_p,n),
                          K(C_{p^\infty},n):p\text{ prime}\}.
 \tag{D.3.7b}
\]
The same conclusion with the universal list $\mathcal E$ holds for
any class of connected polyGEM generators, in particular for any
class of connected GEMs or nilpotent finite Postnikov pieces.
Consequently these localizers
are accessible, each has a countable CW complex as a single reversible
aspherical generator, and there are at most $2^{\aleph_0}$ of them.
These conclusions hold in the fixed small universe and use no large
cardinal hypothesis.

More precisely, if $A\ne0$ is torsion-free and
$J=\{p:pA=A\}$, then
\[
 \min\mathrm{Rev}(K(A,n))
    =\min\mathrm{Rev}(K(\mathbb Z[J^{-1}],n)).
 \tag{D.3.7c}
\]
For a nonzero $p$-primary torsion group $A$, its ordinary nullity,
and hence its reversible closure, is that of $K(C_{p^\infty},n)$
if $A$ is divisible, and that of $K(\mathbb F_p,n)$ otherwise.
The list is a list of sufficient normal forms; no assertion that all
its subsets generate distinct localizers is intended.
\end{claim}

\begin{proof}
We write $\mathcal Q=R(\loc(H\mathbb Q))$ and write
$\operatorname{Rev}(X)$ for the smallest reversible localizer with
$X$ aspherical.  Ordinary $X$-acyclicity implies asphericity for
$\operatorname{Rev}(X)$.  We use ordinary nullity only where explicitly
stated.  Its acyclic objects are closed under retracts, pointed
homotopy colimits and fibre extensions
\citep[beginning of Section~5]{EH5Bousfield1997}; in particular they
are closed under finite products.

\medskip\noindent
\emph{Rational changes of abelian groups.}
First, $\operatorname{Rev}(K(\mathbb Q,n))=\mathcal Q$ for every
$n\geq1$: its suspension spectrum is nonzero and rational, so D.2.5
applies to its reversible closure.  If $A\otimes\mathbb Q\ne0$,
the telescope of multiplication by $2,3,4,\ldots$ on $K(A,n)$ is
$K(A\otimes\mathbb Q,n)$.  The latter has $K(\mathbb Q,n)$ as a
retract.  Thus
\[
 A\otimes\mathbb Q\ne0\quad\Longrightarrow\quad
                \mathcal Q\subseteq\operatorname{Rev}(K(A,n)).
 \tag{D.3.7d}
\]

If $0\to U\to V\to D\to0$ is exact and $D$ is a rational vector
space, then $K(U,n)\to K(V,n)$ is an $H\mathbb F_\ell$-equivalence
for every prime $\ell$.  Indeed, in
\[
 K(U,n)\longrightarrow K(V,n)\longrightarrow K(D,n)
\]
the base is mod-$\ell$ acyclic and its action on the homology of the
fibre is trivial.  For $n=1$ the latter assertion follows from
commutativity of the abelian groups; for $n\geq2$ the base is simply
connected.  The Serre spectral sequence proves the assertion also at
$n=1$, without a claim that a disconnected homotopy fibre is
aspherical.  A rational kernel instead gives a rational fibre and the
same homology conclusion.  A map between connected spaces which is a
mod-$\ell$ homology equivalence for every $\ell$ has a homotopy
cofiber with uniquely divisible integral homology, by the integral
universal coefficient sequences.  Its suspension spectrum is
rational: the fibre of its rationalization is bounded below and
integrally homology-acyclic, and hence zero by the stable Hurewicz
theorem.  Such a map is therefore in $\mathcal Q$.

Consequently kernels or cokernels which are rational vector spaces
may be discarded when comparing reversible aspherical classes that
already contain $\mathcal Q$.  All the comparisons below involving
large products use this fact, not closure under arbitrary products.

\medskip\noindent
\emph{Torsion groups.}
Let $0\ne A$ be $p$-primary.  If $A$ is divisible, the elementary
structure theorem for divisible groups expresses it as a nonempty
direct sum of copies of $C_{p^\infty}$.  Direct-summand retracts and
filtered colimits of finite products give equality of ordinary
nullity with $K(C_{p^\infty},n)$.

If $A$ is not divisible, then $A/pA\ne0$, since all other primes
already act invertibly on $A$.  Thus
\[
 [K(A,n),K(\mathbb F_p,n)]_*
   =H^n(K(A,n);\mathbb F_p)
   =\operatorname{Hom}(A,\mathbb F_p)\ne0.
\]
The exact dichotomy in
\citet[Lemma~9.13]{EH5Bousfield1997} says that
$K(\mathbb F_p,n)$ is either local or acyclic for any localization.
It is not $K(A,n)$-null, so it is $K(A,n)$-acyclic.  Conversely,
$A$ is the filtered union of its finite subgroups.  A finite abelian
$p$-group has a finite filtration with quotients $\mathbb F_p$.
The associated fibre extensions of Eilenberg--Mac Lane spaces,
followed by the filtered homotopy colimit, make $K(A,n)$
$K(\mathbb F_p,n)$-acyclic.  This proof works for every $n\geq1$.

For a general torsion group $A=\bigoplus_p A_p$, each $K(A_p,n)$
is a retract of $K(A,n)$, and the latter is the filtered homotopy
colimit of the finite products of these factors.  This proves the
torsion part of the claim using only the two stated types at each
prime.

\medskip\noindent
\emph{Two module facts, including the required realization.}
We record the algebraic and homotopical details needed for a mixed
group.  Every $\mathbb Z_p$-module $M$, viewed as an abstract abelian
group, satisfies
\[
              K(M,n)\text{ is }K(\mathbb Z_p,n)\text{-acyclic}.
 \tag{D.3.7e}
\]
The discrete valuation ring $\mathbb Z_p$ is a principal ideal
domain, so $M$ has a free resolution
$0\to F_1\to F_0\to M\to0$ of length one.  Each $K(F_i,n)$ is
a filtered homotopy colimit of finite products of
$K(\mathbb Z_p,n)$.  In simplicial abelian group models, realize the
bar construction for the additive action of $K(F_1,n)$ on
$K(F_0,n)$.  Its terms are
$K(F_0,n)\times K(F_1,n)^j$ and are all
$K(\mathbb Z_p,n)$-acyclic.  Its realization is $K(M,n)$: the
homotopy quotient has fibre sequence
\[
 K(F_0,n)\longrightarrow
 K(F_0,n)\mathbin{\!//\!}K(F_1,n)
                  \longrightarrow K(F_1,n+1),
\]
and the homotopy exact sequence, with injective connecting map
$F_1\to F_0$, leaves just $M$ in degree $n$.  Homotopy-colimit
closure proves (D.3.7e).  This is a genuine bar quotient; it does not
use a two-out-of-three rule for arbitrary homotopy fibres.

Next let $F$ be a torsion-free abelian group, uniquely divisible by
every prime other than $p$.  Its $p$-adic completion
\[
                 \widehat F_p=\lim_k F/p^kF
\]
is a torsion-free, separated and complete $\mathbb Z_p$-module, and
$\widehat F_p/p\widehat F_p=F/pF$.  For completeness, torsion-freeness
of $F$ makes division by $p$ unique modulo $p^k$ whenever an element
vanishes modulo $p$.  Applied to compatible residue classes, this
identifies the kernel of
$\widehat F_p\to F/pF$ with $p\widehat F_p$ and shows that
$p$-torsion in the limit is zero.  The same argument with $p^k$
gives all the displayed completeness assertions.

Every torsion-free separated complete $\mathbb Z_p$-module $M$ is
the completion of a free $\mathbb Z_p$-module.  Choose lifts of a
basis of $M/pM$ and map the completion of the free module on this
basis to $M$.  Successive finite expressions modulo $p$ give
surjectivity by $p$-adic convergence.  An element of the kernel is
divisible by $p$ in the source; torsion-freeness of the target puts
its quotient back in the kernel.  Iteration and separatedness make
the kernel zero.  If $M/pM\ne0$, a chosen coordinate of this
completed free module exhibits $\mathbb Z_p$ as a retract of $M$.
The completed free module need not be free as an abstract module.
Only the preceding length-one resolution uses abstract freeness.

\medskip\noindent
\emph{The completion groups of a reduced abelian group.}
Split $A=D\oplus A_0$, with $D$ its maximal divisible subgroup and
$A_0$ reduced.  Divisible groups are injective, and
\[
 D\cong\mathbb Q^{(\kappa_0)}
              \oplus\bigoplus_p C_{p^\infty}^{(\kappa_p)}.
\]
Its contribution is already on the stated list.  We may suppose
$A\otimes\mathbb Q\ne0$, since torsion groups were dealt with
above, and hence retain $\mathcal Q$ among the reversible classes
being compared.  Put
\[
                   C_p=\operatorname{Ext}^1_{\mathbb Z}
                                  (C_{p^\infty},A_0).
\]
Applying $\operatorname{Hom}(-,A_0)$ to
$0\to\mathbb Z\to\mathbb Q\to\mathbb Q/\mathbb Z\to0$
gives
\[
 0\longrightarrow A_0\longrightarrow\prod_p C_p
       \longrightarrow\operatorname{Ext}^1_{\mathbb Z}
                                 (\mathbb Q,A_0)\longrightarrow0.
 \tag{D.3.7f}
\]
Here $\operatorname{Hom}(\mathbb Q,A_0)=0$ because the image of
a divisible group is divisible, and the final Ext group is a
rational vector space by the action of $\mathbb Q$ on its first
argument.  Thus (D.3.7f) identifies $K(A_0,n)$ and
$K(\prod_p C_p,n)$ modulo $\mathcal Q$.

Each $C_p$ is a $\mathbb Z_p$-module by the endomorphism action on
$C_{p^\infty}$.  In particular, multiplication by any prime
$q\ne p$ is invertible.  Also
\[
       \operatorname{Hom}(\mathbb Q,C_p)=0,
       \qquad \operatorname{Ext}^1(\mathbb Q,C_p)=0.
 \tag{D.3.7g}
\]
Here is a direct justification rather than an assumption of ordinary
separated completeness.  Since $A_0$ is reduced,
$\operatorname{Hom}(C_{p^\infty},A_0)=0$.  Since $\mathbb Z$ has
global dimension one, the derived group
$\mathbf R\operatorname{Hom}(C_{p^\infty},A_0)$ has its only
cohomology $C_p$ in degree one.  Derived tensor--Hom adjunction gives
\[
 \mathbf R\operatorname{Hom}\bigl(\mathbb Q,
       \mathbf R\operatorname{Hom}(C_{p^\infty},A_0)\bigr)
  \simeq
 \mathbf R\operatorname{Hom}
       (\mathbb Q\otimes^{\mathbf L}C_{p^\infty},A_0)=0,
\]
which proves (D.3.7g).  This is the Ext-completeness condition of
\citet[Section~5.6]{EH5Bousfield1997}.  It does not assert that the
possibly torsion group $C_p$ is separated for its ordinary $p$-adic
topology.

The quotient
\[
                     (\prod_p C_p)/(\bigoplus_p C_p)
 \tag{D.3.7h}
\]
is a rational vector space.  Indeed, for each prime $\ell$, the
kernel and cokernel of multiplication by $\ell$ on either the sum
or the product are exactly those on its $\ell$-coordinate.  The
maps on these kernels and cokernels induced by the inclusion are
isomorphisms.  The snake lemma makes multiplication by $\ell$ an
isomorphism on (D.3.7h).  Hence modulo $\mathcal Q$ the contribution
of the product is the join of the individual $K(C_p,n)$
contributions: the sum is a filtered homotopy colimit of finite
products, and each coordinate is a retract of the product.

\medskip\noindent
\emph{Reduction at one prime.}
Fix $p$, set $M=C_p$, let $T$ be its torsion subgroup, and put
$F=M/T$.  The torsion in $M$ is $p$-primary.  Both $M$ and $F$ are
uniquely divisible by all other primes, and $F$ is torsion-free.
The long Ext sequence and (D.3.7g) give
$\operatorname{Ext}^1(\mathbb Q,F)=0$, since
$\operatorname{Ext}^2_{\mathbb Z}=0$.

If $F/pF=0$, then $F$ is a rational vector space.  Thus $K(T,n)$
and $K(M,n)$ agree modulo $\mathcal Q$.  The group $T$ is reduced:
a nonzero divisible subgroup would contain a copy of
$C_{p^\infty}$ and receive a nonzero map from $\mathbb Q$,
contradicting (D.3.7g).  When $T\ne0$ its ordinary nullity is
therefore that of $K(\mathbb F_p,n)$ by the torsion argument.
When $T=0$ the group $M=F$ is rational and (D.3.7g) makes $M=0$.

Suppose instead that $F/pF\ne0$.  The same exact sequence used in
(D.3.7f), now applied to the torsion-free group $F$, gives
\[
 0\longrightarrow D_F\longrightarrow F
                    \longrightarrow\widehat F_p\longrightarrow0,
 \tag{D.3.7i}
\]
where $D_F$ is a rational vector space.  The other prime completions
vanish, and the cokernel is
$\operatorname{Ext}^1(\mathbb Q,F)=0$.  We used here the precise
completion formula
\[
 \operatorname{Ext}^1(C_{p^\infty},F)=\lim_k F/p^kF
\]
for torsion-free $F$; more generally it holds for bounded
$p$-torsion, as stated in
\citet[Section~4, p.~41]{Bousfield1996ULP}.  The subgroup $D_F$
is injective, so (D.3.7i) splits.  The completed-free-module
argument above supplies a surjection
$\widehat F_p\to\mathbb Z_p$.  Its composite gives a surjection
$h:M\to\mathbb Z_p$.

Every additive map $h:M\to\mathbb Z_p$ is
$\mathbb Z_p$-linear.  For $r\in\mathbb Z_p$, choose integers
$a_k\equiv r\pmod {p^k}$; then
$h(rx)-a_kh(x)\in p^k\mathbb Z_p$ for all $k$.
Separatedness of the target gives $h(rx)=rh(x)$.  Choosing
$x\in M$ with $h(x)=1$ therefore splits $h$ by $r\mapsto rx$.
Thus $K(\mathbb Z_p,n)$ is a retract of $K(M,n)$, while (D.3.7e)
gives the reverse ordinary acyclicity.  In this case $K(M,n)$ and
$K(\mathbb Z_p,n)$ have the same ordinary nullity.

Finally $\mathbb Z_p/\mathbb Z_{(p)}$ is a rational vector space.
Multiplication by $p$ is injective on the quotient because
$\mathbb Q\cap\mathbb Z_p=\mathbb Z_{(p)}$, and is surjective by
subtracting an integer with the prescribed residue modulo $p$;
every other prime is already invertible on both groups.  Consequently
$K(\mathbb Z_p,n)$ can be replaced by the countable
$K(\mathbb Z_{(p)},n)$ in the presence of $\mathcal Q$.
This completes the reduction of every mixed abelian group to
(D.3.7b).

\medskip\noindent
\emph{The torsion-free formula, arbitrary GEMs and polyGEMs.}
For a torsion-free group $A$, the argument before reduction at one
prime has
$C_p=\widehat A_p$; it is zero exactly when $A/pA=0$.
Every nonzero $C_p$ is torsion-free and has a $\mathbb Z_p$
retract, so its contribution is precisely the
$K(\mathbb Z_{(p)},n)$ contribution modulo $\mathcal Q$.
This profile is identical for $A$ and
$\mathbb Z[J^{-1}]$, proving (D.3.7c).

A connected GEM $G=\prod_{n\geq1}K(A_n,n)$ is weakly equivalent
to the homotopy colimit of its finite partial products: at any fixed
homotopy degree all sufficiently late new factors are more highly
connected.  Each factor is a retract.  Thus its reversible closure
is the join of the closures of its Eilenberg--Mac Lane factors.
For any class of such GEM generators, apply (D.3.7b) to all these
factors and take the union of the resulting subsets of the single
countable list $\mathcal E$.  The wedge of that subset is a
countable CW complex: each listed abelian group has a countable
simplicial Eilenberg--Mac Lane model.  Its reversible closure is
the required localizer, accessible by B.4.  There are only
$2^{\aleph_0}$ subsets of the list, proving the cardinal bound.

The extension to polyGEMs uses a precise ordinary-acyclicity theorem,
not closure under inverse limits.  In the terminology of
\citet{FGHCFFS2015}, a GEM is a $1$-polyGEM, and an $s$-polyGEM
is a homotopy retract of the fibre of a map from an
$(s-1)$-polyGEM to a GEM.  Their Proposition~6.2(4) identifies the ordinary
acyclic classes generated by a polyGEM $X$ and by its free abelian
group construction $\mathbb ZX$.  For connected $X$, Dold--Thom
identifies $\mathbb ZX$ with the GEM
$\prod_{n\geq1}K(\widetilde H_n(X;\mathbb Z),n)$.
Mutual ordinary acyclicity gives equality of their reversible
closures, so the preceding countable-list reduction applies to
every class of connected polyGEMs.  Finally,
\citep[Proposition~5.2]{FGHCFFS2015} makes every connected
nilpotent finite Postnikov piece a polyGEM.  These two cited
statements impose no cardinality or finite-generation bound.
\end{proof}

The ordinary nullity classes obtained from rigid abelian groups in
\citet[Theorems~8.2 and~8.5]{FGHStanley2006} therefore cannot supply
arbitrarily many reversible classes.  Their unsuspended
$K(A,n)$ separators do not survive this reduction.  D.3.7 proves a
set-size statement for all classes generated by connected polyGEMs,
including their intersection with any specified stable-part fibre.
It does not assert that an arbitrary reversible localizer is
generated by polyGEMs, and does not settle the general set-size question
in D.3.2.

\medskip\noindent
\emph{Consequence for the projective-space interval.}
Let $p$ be any prime.  A reversible localizer generated by connected
infinite-loop spaces and satisfying
\[
 \operatorname{Rev}(B\mathbb Z/p)\ \subseteq\ W
          \ \subseteq\ R(\loc(\mathbb S/p))
\]
must equal $\operatorname{Rev}(B\mathbb Z/p)$.  Indeed, each
aspherical generator $X$ has trivial reduced
$\mathbb Z[1/p]$-homology.  A connected infinite-loop space is
nilpotent, so the nilpotent homological localization criterion
\citep[Theorems~6.1.1--6.1.2]{FGHMayPonto2012} gives
$\pi_i(X)[1/p]=0$ for every $i\geq1$.  Choose its connected
connective spectrum $E$, so $X\simeq\Omega^\infty E$ and
$E[1/p]=0$.  By \citet[Theorem~3.2]{EH5Levy2025}, $X$ is already
ordinarily $B\mathbb Z/p$-acyclic.  Thus every generator belongs
to the lower endpoint and $W$ equals that endpoint.  The same
conclusion holds for connected polyGEM generators, by replacing
each by the ordinarily equivalent GEM just used in the proof.
Hence such generators cannot produce a strict intermediate in
the interval of D.3.2.
